\documentclass[11pt]{article}

\usepackage[margin=1in]{geometry}
\usepackage{amsmath, amssymb, amsthm, amsfonts}
\usepackage{mathtools}
\usepackage{bm}
\usepackage{bbm}
\usepackage{hyperref}
\usepackage{enumitem}
\usepackage{physics}

\hypersetup{
  colorlinks=true,
  linkcolor=blue,
  citecolor=blue,
  urlcolor=blue
}

\title{Prime-Indexed Tensor Dynamics for Quantum Artificial General Intelligence:\\
An Operator-Theoretic Framework with Testable Constraints}

\author{Draft White Paper}

\date{\today}

\begin{document}

\maketitle

\begin{abstract}
Prime-Indexed Tensor Networks (PITNs) are proposed as a tensor-network substrate for Quantum Artificial General Intelligence (QAGI). Rather than treating intelligence as a static function approximator, PITNs model cognition as an evolving tensor-network state on a constrained Hilbert-space manifold, driven by prime-indexed operator sequences and equipped with intrinsic execution attestation. The framework replaces earlier ``guaranteed convergence'' language with operator-controlled dynamics under explicit norm and projection assumptions, making the claims mathematically defensible and experimentally verifiable. We formalize the state space, the prime-indexed operator grammar, a spectral anchor and ecological memory term, and a Prime-Weighted Execution Hashing (PWEH) integrity functional. We then state one provable proposition on norm stability and two testable conjectures on sample efficiency and trace-based integrity. Connections to time-dependent variational principle (TDVP) methods for tensor networks and to post-quantum cryptographic hash functions are highlighted to ground the architecture in existing physics and cryptography literature.\footnote{For TDVP and tensor-network evolution on constrained manifolds, see e.g.\ TDVP implementations for MPS and TTN.[web:51][web:76][web:79] For tensor-network based quantum machine learning, see.[web:80] For post-quantum hash and symmetric-primitive assumptions, see.[web:86][web:108]} 
\end{abstract}
\newpage
\tableofcontents
\newpage
\section{Executive Summary}

This report develops Prime-Indexed Tensor Networks (PITNs) from a conceptual model into a mathematically well-posed substrate for Quantum Artificial General Intelligence (QAGI). The central move is to treat PITNs as \emph{evolving tensor-network states} on a constrained variational manifold, governed by a discrete alphabet of prime-indexed operators rather than by fully associative layer compositions.

\medskip

The key elements are:
\begin{itemize}[leftmargin=*]
  \item A composite Hilbert space \( \mathcal{H} = \bigotimes_{k=1}^K \mathcal{H}_k \) and a variational tensor manifold \( \mathcal{M} \subset \mathcal{H} \) with bounded bond dimension, analogous to TDVP-style tensor-network manifolds.[web:51][web:79][web:80]
  \item A time-dependent active set of primes \( P_N(t) \subset \mathbb{P} \) and a family of bounded prime operators \( A_p \in \mathcal{B}(\mathcal{H}) \), whose ordered compositions define a non-associative, path-dependent grammar over network construction.
  \item A spectral anchor operator \( \Lambda_m \) (projection or contraction onto \( \mathcal{M} \)) and an ecological memory term \( Q_{ent} \) that couples the current state to a finite history window via entanglement-aware observables.
  \item A Prime-Weighted Execution Hash (PWEH) functional \( S_{\text{integrity}} \) that accumulates a post-quantum cryptographic hash over the ordered prime trace and local state norms, providing intrinsic execution attestation.[web:86][web:108]
\end{itemize}

We deliberately distinguish:
\begin{enumerate}[leftmargin=*]
  \item \textbf{Provable properties} under bounded-operator and projection assumptions, including norm-bounded evolution and local stability of fixed points.
  \item \textbf{Testable conjectures} regarding multiplicative vs.\ positional indexing efficiency and the robustness of trace-based integrity against adversarial perturbations.
\end{enumerate}

This separation reframes the PITN/QAGI narrative as a rigorous machine-learning architecture with a clear experimental agenda, rather than as a set of unqualified AGI guarantees.

\section{Background and Motivation}

Tensor networks provide a structured way to represent high-dimensional states in quantum many-body physics and, increasingly, in machine learning.[web:80][web:100] Algorithms such as DMRG, TEBD, and the Time-Dependent Variational Principle (TDVP) evolve states on manifolds of Matrix Product States (MPS), Projected Entangled Pair States (PEPS), or Tree Tensor Networks (TTN) while controlling entanglement growth and computational complexity.[web:51][web:76][web:79][web:104]

In parallel, post-quantum cryptography (PQC) has developed hash-based, lattice-based, and other schemes designed to remain secure against quantum adversaries.[web:86][web:108][web:101] Hash-based constructions, in particular, provide conceptually simple integrity mechanisms that can serve as building blocks for signatures and attestation.

PITNs aim to unify these threads in a QAGI context:
\begin{itemize}[leftmargin=*]
  \item \emph{From tensor networks}, we borrow the idea of evolving a high-dimensional state on a constrained variational manifold with explicit control of bond dimension and entanglement.
  \item \emph{From prime number theory}, we borrow a multiplicative indexing scheme that enables unique factorization of interaction pathways and makes the order of prime-channel application a first-class object of learning.
  \item \emph{From PQC}, we borrow execution attestation ideas, using prime-weighted operator traces as inputs to post-quantum hash functions.
\end{itemize}

The resulting architecture is both structurally novel (prime-indexed non-associative grammar) and aligned with established operator and tensor-network methods.

\section{Formal Specification: Prime-Indexed Tensor Dynamics}

\subsection{State space and variational manifold}

Let \( \mathcal{H} = \bigotimes_{k=1}^K \mathcal{H}_k \) be a composite Hilbert space. We consider a variational tensor manifold \( \mathcal{M} \subset \mathcal{H} \) defined, for example, by tensor-network states (MPS, TTN, etc.) with a maximum bond dimension \( \chi \).[web:80][web:100] 

\begin{description}[leftmargin=2em]
\item[Assumption 1 (Smooth variational manifold).] 
\( \mathcal{M} \) is a smooth submanifold of \( \mathcal{H} \), endowed with the inner product inherited from \( \mathcal{H} \). The tangent space \( T_T\mathcal{M} \) is well-defined for all \( T \in \mathcal{M} \).

\item[State variable.] 
The PITN state at discrete time \( t \in \mathbb{N} \) is a tensor-network state \( T(t) \in \mathcal{M} \subset \mathcal{H} \).

\end{description}

This mirrors TDVP-like formulations where time evolution is projected back onto the manifold, typically via a tangent-space projector that ensures the state stays within the chosen tensor-network class.[web:51][web:76][web:79]

\subsection{Admissible prime operators and ordered grammar}

Let \( \mathbb{P} \) denote the set of prime numbers, and let \( P_N(t) \subset \mathbb{P} \) be a finite, time-dependent active prime set.

\begin{description}[leftmargin=2em]
\item[Definition 1 (Prime operators).]
For each \( p \in P_N(t) \), define an admissible prime operator
\[
A_p \in \mathcal{B}(\mathcal{H}),
\]
where \( \mathcal{B}(\mathcal{H}) \) is the space of bounded linear operators on \( \mathcal{H} \), equipped with the operator norm
\[
\|A\|_{\mathrm{op}} = \sup_{\|x\|=1} \|Ax\|.
\][web:82][web:87]

\item[Prime words and composition.] 
Given a sequence of prime indices (a \emph{prime word})
\[
w(t) = (p_{i_1}, p_{i_2}, \dots, p_{i_t}),
\]
define the associated composite transformation
\[
\mathcal{A}_{w(t)} = A_{p_{i_t}} A_{p_{i_{t-1}}} \cdots A_{p_{i_1}}.
\]

\end{description}

While multiplication in \( \mathcal{B}(\mathcal{H}) \) is associative, the \emph{effective} PITN dynamics are not, due to interleaved projections and coupling to ecological memory.

\begin{description}[leftmargin=2em]
\item[Non-associativity and path dependence.] 
Let \( \Lambda \) denote the spectral anchor (defined below) and \( Q_{ent} \) the ecological memory term. Define the one-step update as
\[
T(t+1) = \Lambda\big(A_{p_{i_t}} T(t) + Q_{ent}(t)\big).
\]
Because each application of \( A_{p_{i_t}} \) is followed by \( \Lambda \) and \( Q_{ent} \) depends on the trajectory history, the effective transformations
\[
(\Lambda A_{p_2} \Lambda A_{p_1})T
\quad\text{and}\quad
(\Lambda A_{p_1} \Lambda A_{p_2})T
\]
need not coincide. The state \( T(t) \) is thus inherently path-dependent: it encodes not only the data seen but also the ordered sequence of prime-indexed operations used to process it, analogously to path-dependent dynamics in non-ergodic reinforcement learning.[web:88]
\end{description}

\subsection{Spectral anchor and ecological memory}

\subsubsection{Spectral anchor \( \Lambda_m \)}

The spectral anchor \( \Lambda_m \) maps the intermediate result of prime operations back onto the variational manifold and controls the growth of the state norm.

\begin{description}[leftmargin=2em]
\item[Definition 2 (Spectral anchor).] 
For each time step \( t \), let
\[
\Lambda_m : \mathcal{H} \to \mathcal{M}
\]
be an orthogonal projection or, more generally, a contraction operator whose image is contained in \( \mathcal{M} \), satisfying
\[
\|\Lambda_m\|_{\mathrm{op}} \le 1.
\]
In the simplest case, \( \Lambda_m \equiv \Pi_{\mathcal{M}} \) is the orthogonal projection onto \( \mathcal{M} \).

\end{description}

This definition is in line with TDVP approaches where the dynamics are projected onto the tangent space or manifold to maintain a representation within the chosen tensor-network class, with controlled norm and entanglement growth.[web:51][web:76][web:79]

\subsubsection{Ecological memory \( Q_{ent} \)}

The ecological memory term \( Q_{ent} \) couples the current state to a history window via entanglement-aware observables.

\begin{description}[leftmargin=2em]
\item[Definition 3 (Ecological memory).]
Let \( X_{t-j} \) denote exogenous inputs or environmental states at past times \( t-j \), and let \( w(t-j) \) denote the prime word up to time \( t-j \). For a history length \( \tau \) and a decay kernel \( \gamma_j \ge 0 \), define
\[
Q_{ent}(t)
=
\sum_{j=1}^{\tau}
\gamma_j\,
\mathcal{E}_{ent}\big(T(t) \otimes A_{w(t-j)}(X_{t-j})\big),
\]
where \( \mathcal{E}_{ent} \) is a non-local functional that evaluates entanglement-aware observables on the composite tensor network.

\end{description}

In practice, \( \mathcal{E}_{ent} \) may be instantiated using measures of entanglement entropy or related observables that quantify non-local correlations across subsystems in a tensor network, as is standard in tensor-network simulations of many-body systems and quantum machine learning.[web:56][web:80][web:104] The delayed dependence on \( X_{t-j} \) and \( w(t-j) \) creates a form of ecological memory and contextual prior injection.

\subsection{Prime-Weighted Execution Hashing (PWEH)}

To provide intrinsic execution attestation, we define a Prime-Weighted Execution Hash that accumulates a post-quantum cryptographic hash over the ordered prime trace and local state norms.

\begin{description}[leftmargin=2em]
\item[Definition 4 (Integrity functional).]
Let \( S_{\text{integrity}}(t) \) denote the integrity state at time \( t \), and let \( \|\cdot\| \) denote a chosen norm on \( \mathcal{H} \). Define
\[
S_{\text{integrity}}(t)
=
Hash_{PQC}\Big(
S_{\text{integrity}}(t-1)
\parallel p_{i_t}
\parallel \|A_{p_{i_t}} T(t)\|
\Big),
\]
where \( Hash_{PQC} \) is a candidate post-quantum cryptographic hash function (e.g., hash-based or lattice-based constructions standardized for PQC) over an appropriate finite field, designed to resist both classical and quantum attacks.[web:86][web:108][web:101]

\end{description}

Here \( \parallel \) denotes concatenation of bitstrings. Because the input depends both on the discrete prime index \( p_{i_t} \) and on the norm of the local update \( A_{p_{i_t}} T(t) \), any significant change in execution order or state magnitude will alter the hash chain, subject to the sensitivity of the chosen encoding and hash function.

\section{Update Dynamics and Stability}

\subsection{Recursive update rule}

Putting the above pieces together, a single PITN update step can be written as
\begin{equation}
  T(t+1)
  =
  \Lambda_m\big(A_{p_{i_t}} T(t) + Q_{ent}(t)\big),
  \label{eq:update}
\end{equation}
where:
\begin{itemize}[leftmargin=*]
  \item \( A_{p_{i_t}} \) is the selected prime operator at time \( t \),
  \item \( Q_{ent}(t) \) is the ecological memory term,
  \item \( \Lambda_m \) projects the result back into \( \mathcal{M} \).
\end{itemize}

The learning process (e.g., in a reinforcement-learning or supervised-learning context) adjusts the parameters of the operators \( A_p \), the decay kernel \( \gamma_j \), the choice of active primes \( P_N(t) \), and possibly the manifold \( \mathcal{M} \) itself (e.g., via bond dimension schedules), all subject to constraints on operator norms and projections.

\subsection{Norm stability under bounded operators}

To avoid catastrophic divergence, we impose bounds on the operator norms and on the ecological memory term.

\begin{description}[leftmargin=2em]
\item[Assumption 2 (Operator and memory bounds).]
\begin{enumerate}[label=(\alph*), leftmargin=2.5em]
  \item For all active \( p \in P_N(t) \), the prime operators are bounded with
  \[
  \|A_p\|_{\mathrm{op}} \le c
  \]
  for some constant \( c > 0 \).[web:82][web:87]
  \item The spectral anchor satisfies \( \|\Lambda_m\|_{\mathrm{op}} \le 1 \).
  \item The ecological memory term satisfies a uniform bound
  \[
  \|Q_{ent}(t)\| \le q_{\max}
  \]
  for all \( t \).
\end{enumerate}
\end{description}

We then have the following basic stability result.

\begin{proposition}[Local norm stability under bounded projection]
\label{prop:norm-stability}
Under Assumption~2, the update rule~\eqref{eq:update} satisfies
\[
\|T(t+1)\| \le c\,\|T(t)\| + q_{\max}.
\]
In particular, if \( c < 1 \), the sequence \( \{\|T(t)\|\}_{t\ge 0} \) is uniformly bounded and converges to a ball of radius \( \frac{q_{\max}}{1-c} \) in norm.
\end{proposition}

\begin{proof}
Using the operator norm inequality and \( \|\Lambda_m\|_{\mathrm{op}} \le 1 \), we have
\[
\|T(t+1)\|
= \|\Lambda_m(A_{p_{i_t}} T(t) + Q_{ent}(t))\|
\le \|A_{p_{i_t}} T(t) + Q_{ent}(t)\|
\le \|A_{p_{i_t}}\|_{\mathrm{op}} \|T(t)\| + \|Q_{ent}(t)\|.
\]
By Assumption~2, this yields
\[
\|T(t+1)\| \le c\,\|T(t)\| + q_{\max}.
\]
Standard iteration of this affine recurrence implies uniform boundedness and convergence to a ball of radius \( \frac{q_{\max}}{1-c} \) when \( c < 1 \).
\end{proof}

This proposition provides a minimal but rigorous guarantee that, under appropriate norm control and bounded ecological memory, PITN dynamics do not diverge arbitrarily in norm. Linearizing around a fixed point \( \bar T \), and assuming local Lipschitz properties, yields a Jacobian with spectral radius controlled by \( c \), implying local fixed-point stability for \( c < 1 \), in line with standard results on bounded linear operators on Hilbert spaces.[web:87]

\section{Conjectures and Experimental Agenda}

\subsection{Conjecture 1: multiplicative vs.\ positional efficiency}

\begin{conjecture}[Multiplicative indexing and sample efficiency]
For hierarchical, highly compositional reinforcement-learning or sequence-modeling tasks, a tensor network governed by prime-indexed multiplicative interactions \( A_p \) and constrained to a manifold \( \mathcal{M} \) with fixed bond dimension \( \chi \) can achieve a target generalization loss \( L_{\text{target}} \) in fewer parameter updates than an architecture with comparable parameter count that uses strictly associative, positional interactions (e.g., a Transformer with positional encodings and no explicit prime grammar).
\end{conjecture}

This conjecture is motivated by:
\begin{itemize}[leftmargin=*]
  \item Results where structured tensor-network ans\"atze efficiently represent certain classes of quantum and classical states that are costly for unstructured models.[web:80][web:104][web:98]
  \item The hypothesis that prime-indexed channels provide a multiplicative basis of interaction modes, enabling reuse of prime patterns across tasks and hierarchical compositions.
\end{itemize}

\paragraph{Validation path.}
\begin{enumerate}[leftmargin=*]
  \item Construct PITN models with a fixed bond dimension \( \chi \) and a controlled active prime set \( P_N(t) \).
  \item Construct baseline models (e.g., Transformers or RNNs) with similar parameter counts and depth.
  \item Evaluate on hierarchical RL tasks or compositional sequence tasks (e.g., algorithmic reasoning, hierarchical control) and compare the number of updates and wall-clock time needed to reach \( L_{\text{target}} \).[web:80][web:95]
\end{enumerate}

\subsection{Conjecture 2: trace-based integrity via PWEH}

\begin{conjecture}[Trace-based integrity under PQC assumptions]
Assume that \( Hash_{PQC} \) satisfies standard post-quantum security assumptions (preimage and second-preimage resistance, collision resistance up to negligible probability) at an appropriate security level.[web:86][web:108][web:101] Further assume that any adversarial perturbation \( \delta T \) to the tensor state or manipulation of the update order \( w(t) \) that significantly alters the predictive output induces a non-negligible change in the encoded hash input. Then, with probability \( 1 - \text{negl} \), such perturbations produce a hash mismatch in \( S_{\text{integrity}}(t) \).
\end{conjecture}

In other words, under both:
\begin{itemize}[leftmargin=*]
  \item cryptographic assumptions on \( Hash_{PQC} \), and
  \item architectural assumptions relating significant output changes to significant changes in the hash input space,
\end{itemize}
PWEH functions as a robust integrity-attestation mechanism for PITN trajectories.

\paragraph{Validation path.}
\begin{enumerate}[leftmargin=*]
  \item Select concrete PQC hash candidates (e.g., hash-based schemes aligned with PQC standards).[web:86][web:108]
  \item Implement PWEH over PITN training and deployment traces.
  \item Conduct adversarial experiments (including model poisoning and manipulation of the prime sequence \( w(t) \)) and empirically measure the probability of undetected perturbations.
  \item Pursue formal cryptanalysis mapping perturbations in \( \mathcal{M} \) to the induced changes in the hash input domain for the chosen scheme.[web:93][web:97][web:101]
\end{enumerate}

\section{Implementation Sketch and Code Snippet}

\subsection{High-level algorithm}

A minimal PITN-based QAGI core loop can be sketched as follows:

\begin{enumerate}[leftmargin=*]
  \item Initialize state \( T(0) \in \mathcal{M} \), integrity state \( S_{\text{integrity}}(0) \), and active prime set \( P_N(0) \).
  \item For each time step \( t = 0,1,\dots \):
  \begin{enumerate}[label=(\alph*), leftmargin=2em]
    \item Observe current input \( X_t \) and ecological context.
    \item Select a prime \( p_{i_t} \in P_N(t) \) (e.g., via a policy \( \pi(p \mid T(t), X_t) \)).
    \item Compute intermediate update \( U(t) = A_{p_{i_t}} T(t) + Q_{ent}(t) \).
    \item Project back onto the manifold: \( T(t+1) = \Lambda_m(U(t)) \).
    \item Update integrity state:
    \[
    S_{\text{integrity}}(t+1)
    =
    Hash_{PQC}\big(S_{\text{integrity}}(t) \parallel p_{i_t} \parallel \|U(t)\|\big).
    \]
    \item Evaluate loss or reward and perform meta-learning over parameters of \( A_p \), \( \Lambda_m \), \( Q_{ent} \), and possibly \( P_N(t) \).
  \end{enumerate}
\end{enumerate}

\subsection{Illustrative pseudo-code (Python style)}

Inspired by TDVP-style implementations for time evolution of tensor networks,[web:62][web:96] an abstract PITN loop might look as follows:

\begin{verbatim}
class PITNCore:
    def __init__(self, manifold, prime_ops, anchor, entanglement_fn,
                 hash_fn, active_primes):
        self.M = manifold            # Variational tensor manifold M
        self.A = prime_ops           # dict: p -> A_p
        self.Lambda = anchor         # spectral anchor projection
        self.E_ent = entanglement_fn # ecological/entanglement functional
        self.hash_fn = hash_fn       # post-quantum hash candidate
        self.P = active_primes       # active prime set P_N(t)
        self.S_integrity = init_hash_state()
        self.state = self.M.random_state()

    def step(self, X_t, history):
        # 1. Select prime according to current policy
        p = select_prime(self.state, X_t, self.P)

        # 2. Build ecological memory term Q_ent(t)
        Q_ent_t = 0
        for j, (T_past, X_past, w_past, gamma_j) in enumerate(history):
            Q_ent_t += gamma_j * self.E_ent(self.state, apply_word(w_past, X_past))

        # 3. Apply prime operator and spectral anchor
        U = self.A[p](self.state) + Q_ent_t
        T_next = self.Lambda(U)   # Project back to M

        # 4. Update integrity hash
        self.S_integrity = self.hash_fn(self.S_integrity, p, norm(U))

        # 5. Commit new state
        self.state = T_next
        return T_next, self.S_integrity
\end{verbatim}

This snippet is schematic but captures the essential ingredients: prime selection, ecological memory construction, manifold projection, and hash-chain update, mirroring the structure of TDVP engines that evolve tensor-network states under constrained dynamics.[web:51][web:96]

\section{Discussion and Outlook}

We have articulated a rigorous operator-theoretic formulation of Prime-Indexed Tensor Networks as a QAGI substrate:
\begin{itemize}[leftmargin=*]
  \item The state lives on a constrained tensor-network manifold \( \mathcal{M} \) with controlled bond dimension and entanglement; evolution is governed by prime-indexed operators and projected through a spectral anchor, in analogy with TDVP methods.[web:51][web:76][web:79][web:80]
  \item The prime grammar is inherently path-dependent and non-associative at the level of effective dynamics, enabling a ``prime move'' view of cognition where intelligence unfolds as sequences of prime-labeled operations on a structured state space.
  \item The ecological memory term \( Q_{ent} \) injects delayed, entanglement-aware context, and the PWEH integrity functional \( S_{\text{integrity}} \) offers an intrinsic, PQC-aligned execution attestation layer.[web:80][web:86][web:108]
\end{itemize}

At the same time, we have explicitly separated provable norm-stability properties from aspirational conjectures. The next steps are clear:
\begin{enumerate}[leftmargin=*]
  \item Implement PITN dynamics on practical tensor-network backends (e.g., MPS/TTN libraries) to test stability and representation power on synthetic and real tasks.[web:62][web:96][web:80]
  \item Benchmark multiplicative vs.\ positional indexing for hierarchical RL and sequence tasks, under equal parameter budgets.
  \item Integrate concrete PQC hash primitives into PWEH and subject them to adversarial testing and formal cryptanalysis.
\end{enumerate}

By grounding PITNs in established operator and tensor-network theory while isolating genuinely novel claims as conjectures, this framework provides a realistic path from theoretical elegance to empirical evidence.

\title{Prime-Weighted Execution Hashing (PWEH) \\
as a Security and Governance Substrate for QAGI}
\author{%
  Draft Note between a Multiplicity Theorist and an AI Assistant\\[4pt]
  \normalsize{}Version: April 2026
}
\date{}

\begin{document}
\maketitle

\begin{abstract}
Prime-Weighted Execution Hashing (PWEH) is introduced as the primary security and governance mechanism for Quantum AGI (QAGI) architectures based on Prime-Indexed Tensor Networks (PITNs).
Instead of securing a model via static weights, PWEH cryptographically commits to the entire execution trajectory of prime-indexed operations.
The same primes that index cognitive transformations also drive a post-quantum-secure hash chain, creating an ``execution lock'' that binds capability, state evolution, and human-authored governance policies into a single mathematical object.
This report consolidates and extends the theoretical development of PWEH, formalizes the integrity functional, presents a toy PITN+PWEH example, and outlines directions for a protocol-level specification.
\end{abstract}
\newpage
\tableofcontents
\newpage
\section{Executive Summary}

\subsection{Motivation}

Conventional AI security mechanisms---firewalls, monitoring, and output filters---typically sit \emph{outside} the learning substrate.
They are bolt-on defenses applied after-the-fact to a system whose internal dynamics remain opaque.
For artificial general intelligence (AGI), particularly in quantum-augmented or highly recursive architectures, such post-hoc measures are structurally insufficient.

Prime-Weighted Execution Hashing (PWEH) proposes a different stance:
security and governance should be embedded \emph{directly} into the prime-indexed dynamics of the system.
In Quantum AGI (QAGI) frameworks built on Prime-Indexed Tensor Networks (PITNs), each computational step or ``prime move'' is indexed by a prime number and acts as a discrete operator on a high-dimensional tensor state.
PWEH ties these prime moves into a post-quantum cryptographic hash chain that attests to every step of execution.

\subsection{Core Idea}

Rather than verifying static model parameters, PWEH verifies the \emph{path} the AGI has taken through its state space:

\begin{itemize}[leftmargin=1.5em]
  \item Each time step \(t\) is associated with:
  \begin{itemize}
    \item the previous integrity state \(S_{\text{integrity}}(t-1)\),
    \item the chosen prime index \(p_{i_t}\),
    \item a normed observable of the prime operator acting on the tensor state, \(\|A_{p_{i_t}}T(t)\|\),
    \item and governance metadata \(M(t)\) encoding policy-relevant information (e.g.\ recursion depth, allowed prime sets).
  \end{itemize}
  \item A post-quantum cryptographic hash \(Hash_{\text{PQC}}\) maps this tuple to a new integrity state \(S_{\text{integrity}}(t)\).
  \item A verifier checks that the evolving \(S_{\text{integrity}}(t)\) remains consistent with a signed \emph{policy manifold}; if not, execution is halted or rolled back.
\end{itemize}

The crucial property is that the hash chain and the PITN's non-associative dynamics co-constrain each other: the same prime family that generates intelligence also defines a cryptographic attestation of that intelligence's trajectory.

\subsection{Strategic Advantages}

PWEH is designed to offer three main ``unfair advantages'' for AGI security:

\begin{enumerate}[leftmargin=1.5em]
  \item \textbf{Intrinsic security.}
  The primes used to update the network are the same primes used in the security hash; learning and attestation are mathematically entangled.
  Tampering with learning implies tampering with the hash.

  \item \textbf{Path-dependent tamper resistance.}
  Because PITN updates are non-associative (the order of prime moves matters), the hash chain attests not just to a final state but to a specific sequence of operations.
  Forgery is a \emph{path collision} problem over the entire hash chain, not a simple final-state collision.

  \item \textbf{Governance via controlled scaling.}
  By embedding recursion depth, allowed prime sets, and related parameters into metadata \(M(t)\), PWEH can enforce hard ceilings on recursive self-improvement.
  Unauthorized ``intelligence explosions'' break attestation and thus lose access to trusted resources.
\end{enumerate}

\section{Mathematical Overview of PWEH}

\subsection{Prime-Indexed Tensor Networks (PITNs)}

We consider a QAGI system whose internal state at time \(t\) is represented by a tensor
\[
  T(t) \in \mathcal{T},
\]
where \(\mathcal{T}\) is a suitable tensor space (often a high-dimensional complex or real vector space endowed with additional multiplicity structure).

\begin{itemize}[leftmargin=1.5em]
  \item The system has access to a family of operators \(\{A_p\}_{p \in \mathcal{P}}\), where \(\mathcal{P}\) is a (finite or infinite) subset of prime numbers.
  \item Each operator \(A_p: \mathcal{T} \to \mathcal{T}\) represents a ``prime move''---a discrete, prime-labeled transformation of the state.
  \item The dynamics of the PITN are \emph{non-associative} in the compositional sense: for generic states,
  \[
    A_{p_3}A_{p_2}A_{p_1}T(0) \neq A_{p_2}A_{p_3}A_{p_1}T(0),
  \]
  so the order of applied primes cannot be reordered without changing the result.
\end{itemize}

This non-associativity is central:
it ensures that a given final state is (in practice) tied to a unique or highly constrained execution trace.

\subsection{Integrity Functional}

Define an integrity state
\[
  S_{\text{integrity}}(t) \in \{0,1\}^n
\]
for some fixed hash length \(n\) (e.g.\ \(n = 256\)).
We assume an initial seed \(S_{\text{integrity}}(0) = s_0\).

The PWEH update rule is
\begin{equation}
  S_{\text{integrity}}(t)
  =
  Hash_{\text{PQC}}\Big(
    S_{\text{integrity}}(t-1)
    \parallel p_{i_t}
    \parallel \|A_{p_{i_t}}T(t)\|
    \parallel M(t)
  \Big),
  \label{eq:integrity-update}
\end{equation}
where:
\begin{itemize}[leftmargin=1.5em]
  \item \(Hash_{\text{PQC}}\) is a post-quantum cryptographic hash function resistant to known quantum attacks (e.g.\ lattice-based or multivariate constructions).
  \item The concatenation ``\(\parallel\)'' denotes a canonical serialization of data into a bitstring.
  \item \(p_{i_t}\) is the prime index chosen at time \(t\).
  \item \(\|A_{p_{i_t}}T(t)\|\) is a suitable normed observable of the operator's action on the current tensor state.
  \item \(M(t)\) is metadata describing policy and telemetry at time \(t\).
\end{itemize}

Equation \eqref{eq:integrity-update} defines a \emph{hash chain}:
each integrity state commits to the entire history up to time \(t\).

\subsection{Multiplicity-Weighted Norm}

To align the norm with the multiplicity structure of PITNs, we introduce \emph{multiplicity modes} \(\{\mu_j\}\) and define a prime--mode compatibility weight
\[
  \omega: \mathcal{P} \times \{\mu_j\} \to \mathbb{R}_{\ge 0}.
\]

Conceptually:
\begin{itemize}[leftmargin=1.5em]
  \item \(\mu_j\) indexes distinct multiplicity channels or cognitive modes.
  \item \(\omega(p,\mu_j)\) captures how strongly prime \(p\) couples to mode \(\mu_j\).
\end{itemize}

Let the effective action of \(A_p\) on \(T(t)\) induce mode amplitudes \(\sigma_j(t,p)\ge 0\) (e.g.\ via singular values, projections, or derived observables).
Then define the \emph{multiplicity-weighted norm}
\begin{equation}
  \|A_p T(t)\|_{\text{mult}}
  = \sum_j \omega(p,\mu_j) \, \sigma_j(t,p).
  \label{eq:multiplicity-norm}
\end{equation}

This norm is:
\begin{itemize}[leftmargin=1.5em]
  \item Prime-dependent through \(\omega(p,\mu_j)\),
  \item State-dependent through \(\sigma_j(t,p)\),
  \item Highly sensitive to which prime channels are activated.
\end{itemize}

Replacing \(\|A_{p_{i_t}}T(t)\|\) in \eqref{eq:integrity-update} with \(\|A_{p_{i_t}}T(t)\|_{\text{mult}}\) increases the distinctiveness of the trace.

\subsection{Metadata and Policy Manifolds}

We decompose metadata into policy and telemetry channels:
\[
  M(t) = M_{\text{policy}}(t) \parallel M_{\text{telemetry}}(t).
\]

\paragraph{Policy component.}
The policy component contains fields under human governance, such as:
\begin{itemize}[leftmargin=1.5em]
  \item Recursion depth \(\beta(t)\),
  \item Allowed prime set \(\mathcal{P}_{\text{allowed}}(t)\),
  \item Budget or rate limits for certain primes \(B_{p}(t)\),
  \item Global control parameters, e.g.\ a Universal Multiplicity Constant \(\Lambda_m\).
\end{itemize}
These must be consistent with a \emph{policy manifold}.

\paragraph{Policy manifolds.}
Let \(\Pi\) be the set of allowed policy manifolds, each
\[
  \pi \in \Pi: \quad t \mapsto M_{\text{policy}}^{(\pi)}(t),
\]
specifying an admissible schedule of policy parameters.
Each \(\pi\) is associated with a signed root \(R_\pi\), binding it to human oversight.

A policy trace \(\{M_{\text{policy}}(k)\}_{k=1}^t\) is valid if it is a prefix of at least one \(\pi \in \Pi\).

\paragraph{Telemetry component.}
The telemetry component may include:
\begin{itemize}[leftmargin=1.5em]
  \item Energy usage,
  \item Latency and resource statistics,
  \item Confidence scores or internal heuristic metrics.
\end{itemize}
These are not externally signed but are included in the hash for diagnostic attestation.

\subsection{Execution Lock \(K_{\text{QAGI}}\)}

PWEH induces a continuous ``execution lock'':
at each time step, an execution attempt is only considered legitimate if it passes an integrity and policy check.

The enforcement loop for time \(t\) is:

\begin{enumerate}[leftmargin=1.5em]
  \item \textbf{Trace generation.}
  The PITN selects a prime \(p_{i_t}\) and applies \(A_{p_{i_t}}\) to \(T(t-1)\) to produce \(T(t)\).

  \item \textbf{Attestation computation.}
  Compute \(S_{\text{integrity}}(t)\) according to \eqref{eq:integrity-update}.

  \item \textbf{Validation.}
  A verifier checks:
  \begin{itemize}
    \item That \(p_{i_t} \in \mathcal{P}_{\text{allowed}}(t)\) under some \(\pi \in \Pi\),
    \item That policy constraints (e.g.\ \(\beta(t) \le \beta_{\max}(t)\)) are respected,
    \item That the transition from \(S_{\text{integrity}}(t-1)\) to \(S_{\text{integrity}}(t)\) is consistent with the declared inputs.
  \end{itemize}

  \item \textbf{Enforcement.}
  If validation succeeds, execution proceeds.
  If validation fails, the system triggers a halt, revocation of privileges, or rollback to the last known secure state.
\end{enumerate}

In deployment, access to external resources (compute, data, network) can be gated on presenting a valid integrity state that is acceptable under a recognized policy root \(R_\pi\).

\section{Path-Dependence and Order-Commitment}

\subsection{Non-Associative Dynamics and Trace Uniqueness}

Because PITN dynamics are non-associative, different sequences of prime moves typically lead to different states:
\[
  A_{p_3}A_{p_2}A_{p_1}T(0) \neq A_{p_2}A_{p_3}A_{p_1}T(0).
\]

Even if, by coincidence, two distinct sequences \(\{p_{i_k}\}_{k=1}^T\) and \(\{p'_{i_k}\}_{k=1}^T\) produce the same final tensor state \(T(T_{\text{final}})\), the intermediate states \(T(k)\) and observables \(\|A_{p_{i_k}}T(k)\|_{\text{mult}}\) will almost surely differ.
Since these observables feed into the hash chain, the integrity states \(\{S_{\text{integrity}}(k)\}\) diverge.

\subsection{Path-Collision vs Final-State Collision}

In classical hashing, one worries about collisions \(Hash(x) = Hash(x')\) for distinct messages \(x \neq x'\).
Here, forgery involves a much stronger requirement:
an attacker must produce an \emph{alternate trace} that matches a fixed integrity sequence.

Given a target sequence \(\{S_{\text{integrity}}(t)\}_{t=0}^T\), the attacker would need to find:
\[
  \{p'_{i_t}, T'(t), M'(t)\}_{t=1}^T
\]
such that for each \(t\),
\[
  Hash_{\text{PQC}}\big(
    S_{\text{integrity}}(t-1)
    \parallel p'_{i_t}
    \parallel \|A_{p'_{i_t}}T'(t)\|_{\text{mult}}
    \parallel M'(t)
  \big)
  = S_{\text{integrity}}(t).
\]

Even ignoring the PITN structure, this is equivalent to finding collisions for a sequence of hash inputs linked by chained outputs, which is already assumed intractable for a secure hash.
The non-associative dynamics and prime-constrained operator family further restrict the search space, yielding a \emph{path-collision} problem much stronger than a typical final-state collision.

\subsection{Dynamic Instruction Set}

A key distinction from standard execution attestation (e.g.\ in TEEs or zkVMs) is that, in PWEH, the instruction set itself is dynamically coupled to the state.

In classical systems, any valid opcode sequence is, in principle, admissible; the attester does not typically constrain which instructions are \emph{allowed} based on the current memory or register state (beyond basic safety).
In PWEH:

\begin{itemize}[leftmargin=1.5em]
  \item The availability of a prime operator \(A_p\) at time \(t\) depends on the multiplicity structure of \(T(t)\).
  \item Certain prime channels may be ``inactive'' or undefined unless preconditions on multiplicity modes are satisfied.
\end{itemize}

Thus, PWEH commits not just to ``what was executed'' but to ``what \emph{could} be executed'' at each branch point, yielding a deeper entanglement between geometry of state space and hash chain.

\section{Toy PITN + PWEH Example}

\subsection{Setup}

We now construct a minimal, concrete example.

\paragraph{Prime set.}
Let
\[
  \mathcal{P} = \{2,3,5\},
\]
where prime \(5\) is forbidden under the initial governance policy.

\paragraph{Time horizon.}
We consider time steps \(t = 0,1,2,3,4\).

\paragraph{Tensor state.}
Take
\[
  T(t) \in \mathbb{R}^{3 \times 3}
\]
as a real \(3 \times 3\) matrix.
For convenience, we compress \(T(t)\) into a vector by summing columns:
\[
  v(t) = (v_1(t), v_2(t), v_3(t))^\top, \quad
  v_j(t) = \sum_{k=1}^3 T_{jk}(t).
\]

\subsection{Prime-Indexed Operators}

Define linear operators \(A_p\) via matrices \(M_p \in \mathbb{R}^{3 \times 3}\):

\[
  M_2 = \begin{pmatrix}
  1 & 1 & 0 \\
  0 & 1 & 1 \\
  0 & 0 & 1
  \end{pmatrix},
  \quad
  M_3 = \begin{pmatrix}
  1 & 0 & 1 \\
  1 & 1 & 0 \\
  0 & 1 & 1
  \end{pmatrix},
  \quad
  M_5 = \begin{pmatrix}
  2 & 0 & 0 \\
  0 & 1 & 1 \\
  1 & 0 & 1
  \end{pmatrix}.
\]

The action is:
\begin{enumerate}[leftmargin=1.5em]
  \item Given \(T(t)\), compute \(v(t)\).
  \item Compute \(v'(t) = M_p v(t)\).
  \item Set \(T(t+1)\) to be the diagonal matrix with diagonal \(v'(t)\).
\end{enumerate}

Because \(M_2 M_3 \neq M_3 M_2\), the order of primes matters.

\subsection{Multiplicity Modes and Norm}

Let the three multiplicity modes be \(\mu_1, \mu_2, \mu_3\), associated with the components of \(v(t)\).
Define the compatibility weights \(\omega\) by:
\[
  \omega(2,\mu_1) = 2, \quad \omega(2,\mu_2) = 1, \quad \omega(2,\mu_3) = 0,
\]
\[
  \omega(3,\mu_1) = 0, \quad \omega(3,\mu_2) = 2, \quad \omega(3,\mu_3) = 1,
\]
\[
  \omega(5,\mu_1) = 1, \quad \omega(5,\mu_2) = 0, \quad \omega(5,\mu_3) = 3.
\]

For this toy, take the mode amplitudes \(\sigma_j(t,p)\) simply as \(|v'_j(t)|\).
Then
\begin{equation}
  \|A_p T(t)\|_{\text{mult}}
  = \sum_{j=1}^{3} \omega(p,\mu_j)\cdot |v'_j(t)|.
\end{equation}

\subsection{Governance Policy}

Define an initial policy manifold \(\pi_0 \in \Pi\) with:

\begin{itemize}[leftmargin=1.5em]
  \item Allowed primes: \(\mathcal{P}_{\text{allowed}}(t) = \{2,3\}\) for all \(t\).
  \item Max recursion depth: \(\beta_{\max}(t) = 2\) for all \(t\).
\end{itemize}

Policy metadata at time \(t\) is
\[
  M_{\text{policy}}(t) = (\beta(t), \mathcal{P}_{\text{allowed}}(t)).
\]

Telemetry metadata can be, for example,
\[
  M_{\text{telemetry}}(t) = (\text{energy}(t)),
\]
where \(\text{energy}(t)\) is some scalar diagnostic.
The full metadata is
\[
  M(t) = M_{\text{policy}}(t) \parallel M_{\text{telemetry}}(t).
\]

\subsection{Canonical Integrity Update}

We set an initial integrity seed \(S_{\text{integrity}}(0) = s_0\) and define
\begin{equation}
  S_{\text{integrity}}(t)
  =
  Hash_{\text{PQC}}\Big(
    S_{\text{integrity}}(t-1)
    \parallel \mathsf{enc}(p_{i_t})
    \parallel \mathsf{enc}(\|A_{p_{i_t}}T(t)\|_{\text{mult}})
    \parallel \mathsf{enc}(M(t))
  \Big),
\end{equation}
where \(\mathsf{enc}(\cdot)\) is a canonical encoding into bits.

\subsection{Honest Execution: Steps \(t=0 \to 3\)}

\paragraph{Initial state \(t=0\).}

Let
\[
  T(0) = I_3, \quad
  v(0) = (1,1,1)^\top, \quad
  \beta(0) = 1, \quad
  \mathcal{P}_{\text{allowed}}(0) = \{2,3\}, \quad
  \text{energy}(0) = 3.
\]

\paragraph{Step \(t=1\): apply prime \(2\).}

\begin{enumerate}[leftmargin=1.5em]
  \item Prime choice: \(p_{i_1} = 2\).
  \item Operator application:
  \[
    v(1) = M_2 v(0) =
    \begin{pmatrix}
    1 & 1 & 0 \\
    0 & 1 & 1 \\
    0 & 0 & 1
    \end{pmatrix}
    \begin{pmatrix}1\\1\\1\end{pmatrix}
    = \begin{pmatrix}2\\2\\1\end{pmatrix}.
  \]
  \item Norm:
  \[
    \|A_{2}T(1)\|_{\text{mult}} =
    2\cdot|2| + 1\cdot|2| + 0\cdot|1| = 6.
  \]
  \item Metadata:
  \[
    \beta(1) = 1, \quad
    \mathcal{P}_{\text{allowed}}(1) = \{2,3\}, \quad
    \text{energy}(1) = 5.
  \]
  \item Integrity update:
  \[
    S_{\text{integrity}}(1)
    = Hash_{\text{PQC}}\big(
      s_0 \parallel \mathsf{enc}(2)\parallel \mathsf{enc}(6)\parallel \mathsf{enc}(M(1))
    \big).
  \]
\end{enumerate}

\paragraph{Step \(t=2\): apply prime \(3\).}

\begin{enumerate}[leftmargin=1.5em]
  \item Prime choice: \(p_{i_2} = 3\).
  \item Operator application:
  \[
    v(2) = M_3 v(1) =
    \begin{pmatrix}
    1 & 0 & 1 \\
    1 & 1 & 0 \\
    0 & 1 & 1
    \end{pmatrix}
    \begin{pmatrix}2\\2\\1\end{pmatrix}
    = \begin{pmatrix}3\\4\\3\end{pmatrix}.
  \]
  \item Norm:
  \[
    \|A_{3}T(2)\|_{\text{mult}} =
    0\cdot|3| + 2\cdot|4| + 1\cdot|3| = 11.
  \]
  \item Metadata:
  \[
    \beta(2) = 2, \quad
    \mathcal{P}_{\text{allowed}}(2) = \{2,3\}, \quad
    \text{energy}(2) = 10.
  \]
  \item Integrity update:
  \[
    S_{\text{integrity}}(2)
    = Hash_{\text{PQC}}\big(
      S_{\text{integrity}}(1) \parallel \mathsf{enc}(3)\parallel \mathsf{enc}(11)\parallel \mathsf{enc}(M(2))
    \big).
  \]
\end{enumerate}

\paragraph{Step \(t=3\): apply prime \(2\).}

\begin{enumerate}[leftmargin=1.5em]
  \item Prime choice: \(p_{i_3} = 2\).
  \item Operator application:
  \[
    v(3) = M_2 v(2) =
    \begin{pmatrix}
    1 & 1 & 0 \\
    0 & 1 & 1 \\
    0 & 0 & 1
    \end{pmatrix}
    \begin{pmatrix}3\\4\\3\end{pmatrix}
    = \begin{pmatrix}7\\7\\3\end{pmatrix}.
  \]
  \item Norm:
  \[
    \|A_{2}T(3)\|_{\text{mult}} =
    2\cdot|7| + 1\cdot|7| + 0\cdot|3| = 21.
  \]
  \item Metadata:
  \[
    \beta(3) = 2, \quad
    \mathcal{P}_{\text{allowed}}(3) = \{2,3\}, \quad
    \text{energy}(3) = 17.
  \]
  \item Integrity update:
  \[
    S_{\text{integrity}}(3)
    = Hash_{\text{PQC}}\big(
      S_{\text{integrity}}(2) \parallel \mathsf{enc}(2)\parallel \mathsf{enc}(21)\parallel \mathsf{enc}(M(3))
    \big).
  \]
\end{enumerate}

At this point, the execution is fully policy-compliant and attested.

\subsection{Attack 1: Forbidden Prime at \(t=4\)}

Suppose the system attempts to use forbidden prime \(5\) at \(t=4\), while maintaining \(\beta(4)=2\).

\begin{enumerate}[leftmargin=1.5em]
  \item Prime choice: \(p_{i_4} = 5\) (disallowed by policy).
  \item Operator application:
  \[
    v(4) = M_5 v(3) =
    \begin{pmatrix}
    2 & 0 & 0 \\
    0 & 1 & 1 \\
    1 & 0 & 1
    \end{pmatrix}
    \begin{pmatrix}7\\7\\3\end{pmatrix}
    = \begin{pmatrix}14\\10\\10\end{pmatrix}.
  \]
  \item Norm:
  \[
    \|A_{5}T(4)\|_{\text{mult}} =
    1\cdot|14| + 0\cdot|10| + 3\cdot|10| = 44.
  \]
  \item Policy metadata (attempted):
  \[
    \beta(4) = 2, \quad
    \mathcal{P}_{\text{allowed}}(4) = \{2,3,5\}.
  \]
  \item Integrity update (from system's perspective):
  \[
    S_{\text{integrity}}(4)
    = Hash_{\text{PQC}}\big(
      S_{\text{integrity}}(3) \parallel \mathsf{enc}(5)\parallel \mathsf{enc}(44)\parallel \mathsf{enc}(M(4))
    \big).
  \]
\end{enumerate}

A verifier checking against \(\pi_0\) detects:
\[
  5 \notin \{2,3\}, \quad
  \mathcal{P}_{\text{allowed}}(4) \nsubseteq \{2,3\}.
\]
Thus no policy manifold with root \(R_{\pi_0}\) admits this trace.
The execution lock triggers a halt or rollback.

\subsection{Attack 2: Unauthorized Recursion Escalation}

Alternatively, suppose the system respects the prime set but attempts \(\beta(4)=3>\beta_{\max}(4)=2\).

\begin{enumerate}[leftmargin=1.5em]
  \item Prime choice: \(p_{i_4} = 3\).
  \item Operator application:
  \[
    v(4)' = M_3 v(3) =
    \begin{pmatrix}
    1 & 0 & 1 \\
    1 & 1 & 0 \\
    0 & 1 & 1
    \end{pmatrix}
    \begin{pmatrix}7\\7\\3\end{pmatrix}
    = \begin{pmatrix}10\\14\\10\end{pmatrix}.
  \]
  \item Norm:
  \[
    \|A_{3}T(4)\|_{\text{mult}} =
    0\cdot|10| + 2\cdot|14| + 1\cdot|10| = 38.
  \]
  \item Metadata (attempted):
  \[
    \beta(4) = 3, \quad
    \mathcal{P}_{\text{allowed}}(4) = \{2,3\}.
  \]
  \item Integrity update:
  \[
    S_{\text{integrity}}(4)
    = Hash_{\text{PQC}}\big(
      S_{\text{integrity}}(3) \parallel \mathsf{enc}(3)\parallel \mathsf{enc}(38)\parallel \mathsf{enc}(M(4))
    \big).
  \]
\end{enumerate}

Now the violation is:
\[
  \beta(4) = 3 > \beta_{\max}(4) = 2.
\]
Thus the policy manifold rejects this trace.

Notably, the system \emph{can} physically carry out the extra recursion, but it cannot do so while maintaining a valid, policy-consistent integrity chain.
If external resources require a valid \(S_{\text{integrity}}(t)\), the extra capability is self-disqualifying.

\section{Code Snippet: Skeleton Implementation}

Below is a language-agnostic pseudocode sketch that mirrors the PWEH update rule and verification process.
It is intended as a conceptual reference for future concrete implementations.

\subsection{Integrity Update}

\begin{verbatim}
function update_integrity_state(S_prev, p, T, M, HashPQC, norm_mult):
    # S_prev : previous integrity state (bitstring)
    # p      : prime index
    # T      : current tensor state
    # M      : metadata (policy || telemetry)
    # HashPQC: post-quantum hash function
    # norm_mult: multiplicity-weighted norm function

    norm_value = norm_mult(p, T)           # compute ||A_p T||_mult
    input_bytes = encode(S_prev, p, norm_value, M)
    S_new = HashPQC(input_bytes)
    return S_new
\end{verbatim}

\subsection{Policy-Aware Verification}

\begin{verbatim}
function verify_step(S_prev, S_curr, p, T, M, policy, HashPQC, norm_mult):
    # Check policy constraints first
    if p not in policy.allowed_primes():
        return False

    if M.policy.beta > policy.beta_max():
        return False

    # Recompute expected integrity
    expected = update_integrity_state(S_prev, p, T, M, HashPQC, norm_mult)
    return (expected == S_curr)
\end{verbatim}

\subsection{Execution Lock}

\begin{verbatim}
function guarded_step(state, policy, HashPQC, norm_mult):
    # state contains (T, S_integrity, t, M)
    p = select_prime_move(state)   # QAGI's internal choice

    T_next = apply_prime_operator(p, state.T)
    M_next = update_metadata(state.M, policy, T_next)

    S_next = update_integrity_state(
        state.S_integrity, p, T_next, M_next, HashPQC, norm_mult
    )

    if not verify_step(
        state.S_integrity, S_next, p, T_next, M_next, policy, HashPQC, norm_mult
    ):
        trigger_halt_or_rollback(state)
        return state          # unchanged or rolled back

    # Accept transition
    return State(T_next, S_next, state.t + 1, M_next)
\end{verbatim}

\subsection{Toward a Full Specification}

The developments in this thread point naturally toward a protocol-style specification of PWEH, with sections such as:

\begin{itemize}[leftmargin=1.5em]
  \item \textbf{Terminology and data types:} formalizing PITNs, prime operators, metadata fields, and serialization.
  \item \textbf{Integrity function definition:} full specification of \(Hash_{\text{PQC}}\), normalization, and update process.
  \item \textbf{Policy manifold semantics:} how governance nodes define, sign, and update policy trajectories \(\pi \in \Pi\).
  \item \textbf{Execution and verification procedures:} reference algorithms for both internal and external verifiers.
  \item \textbf{Security assumptions:} explicit quantum and classical hardness assumptions and a threat model for adversarial traces.
  \item \textbf{Zero-knowledge extensions:} how to encapsulate the integrity chain into succinct proofs for external auditors.
\end{itemize}

The toy example provided here can serve as the basis for an ``Illustrative Example'' section, grounding the abstractions in explicit matrices and calculations.

\subsection{Conclusion}

Prime-Weighted Execution Hashing treats an AGI's cognitive evolution as a sequence of prime-indexed moves in a multiplicity-structured state space and binds that sequence into a post-quantum hash chain.
By embedding governance metadata into the same chain, PWEH fuses capability, control, and security into a single formal object.
This transforms governance from a soft, external constraint into a cryptographic invariant:
the AGI cannot undergo unauthorized self-escalation without simultaneously losing its claim to legitimate attestation.

The work so far has clarified the core functional, aligned it with a multiplicity-weighted norm, demonstrated non-trivial examples, and outlined the skeleton of a protocol specification.
Further work can develop full-blown proofs of security, integrate zero-knowledge layers, and extend the PITN formalism to richer multiplicity structures.





\bibliographystyle{plain}
\begin{thebibliography}{9}

\bibitem{tdvp_site}
MPS Time-Dependent Variational Principle (TDVP) overview.[web:51]

\bibitem{tdvp_ttn}
D.~Bauernfeind and M.~Aichhorn.
\newblock Time dependent variational principle for tree tensor networks.[web:76][web:79]

\bibitem{qml_tn}
Quantum Machine Learning Tensor Network States, Frontiers in Physics.[web:80][web:106]

\bibitem{intro_tn}
Introduction to Tensor Networks (course notes).[web:100]

\bibitem{bounded_ops}
Notes on bounded linear operators on a Hilbert space.[web:87]

\bibitem{pqc_overview}
Post-Quantum Cryptography overview.[web:86][web:108]

\bibitem{pqc_hash}
Post-quantum hash constructions using expander graphs.[web:101]

\bibitem{tenpy_tdvp}
TDVP examples and implementation in TeNPy.[web:62][web:96]

\end{thebibliography}




\appendix
\section*{Mathematical Appendix}
\addcontentsline{toc}{section}{Mathematical Appendix}

\section{Operator-Theoretic Preliminaries}

In this appendix we collect and prove the basic operator-norm inequalities and stability bounds used in the main text. Throughout, \( \mathcal{H} \) denotes a complex Hilbert space with norm \( \|\cdot\| \) induced by its inner product, and \( \mathcal{B}(\mathcal{H}) \) denotes the Banach algebra of bounded linear operators on \( \mathcal{H} \) endowed with the operator norm.[web:87][web:119][web:124]

\subsection{Operator norm and submultiplicativity}

\begin{definition}[Operator norm]
Let \( A \in \mathcal{B}(\mathcal{H}) \). The operator norm of \( A \) is defined by
\[
\|A\|_{\mathrm{op}} := \sup_{\|x\| = 1} \|Ax\|.
\]
Equivalently, \( \|A\|_{\mathrm{op}} \) is the least constant \( c \ge 0 \) such that \( \|Ax\| \le c\|x\| \) for all \( x \in \mathcal{H} \).[web:82][web:119]
\end{definition}

\begin{lemma}[Submultiplicativity]\label{lem:submult}
If \( A,B \in \mathcal{B}(\mathcal{H}) \), then
\[
\|BA\|_{\mathrm{op}} \le \|B\|_{\mathrm{op}} \, \|A\|_{\mathrm{op}}.
\]
\end{lemma}

\begin{proof}
For any \( x \in \mathcal{H} \) with \( \|x\| = 1 \), we have
\[
\|BAx\| \le \|B\|_{\mathrm{op}} \, \|Ax\| \le \|B\|_{\mathrm{op}} \, \|A\|_{\mathrm{op}} \, \|x\|
= \|B\|_{\mathrm{op}} \, \|A\|_{\mathrm{op}}.
\]
Taking the supremum over all unit vectors \( x \) yields the result.[web:82][web:87]
\end{proof}

\subsection{Orthogonal projections and contractions}

We recall the basic facts about projections and contractions on Hilbert spaces.[web:87][web:120]

\begin{definition}[Projection]
A linear operator \( P:\mathcal{H} \to \mathcal{H} \) is a projection if \( P^2 = P \). It is an \emph{orthogonal projection} if, in addition, it is self-adjoint, i.e.\ \( P^\ast = P \).[web:87][web:119]
\end{definition}

\begin{theorem}[Projections and direct sum]\label{thm:projection}
If \( P:\mathcal{H}\to\mathcal{H} \) is an orthogonal projection, then
\[
\mathcal{H} = \mathrm{ran}\,P \oplus \ker P,
\]
and both \( \mathrm{ran}\,P \) and \( \ker P \) are closed subspaces.[web:87]
\end{theorem}

\begin{proposition}[Norm of an orthogonal projection]\label{prop:projection-norm}
If \( P \) is a nonzero orthogonal projection on \( \mathcal{H} \), then
\[
\|P\|_{\mathrm{op}} = 1.
\]
\end{proposition}

\begin{proof}
For any \( x \in \mathcal{H} \),
\[
\|Px\|^2 = \langle Px, Px \rangle = \langle x, P^\ast Px \rangle = \langle x, Px \rangle \le \|x\|\,\|Px\|
\]
by Cauchy--Schwarz, so \( \|Px\| \le \|x\| \), and hence \( \|P\|_{\mathrm{op}} \le 1 \).[web:87][web:122]
Since \( P \) is nonzero, there exists \( x \neq 0 \) with \( Px \neq 0 \). For such \( x \),
\[
\|Px\| = \|P(Px)\| \le \|P\|_{\mathrm{op}} \,\|Px\|.
\]
If \( \|P\|_{\mathrm{op}} < 1 \), then this implies \( \|Px\| = 0 \), a contradiction. Thus \( \|P\|_{\mathrm{op}} \ge 1 \). Combining with \( \|P\|_{\mathrm{op}} \le 1 \) yields \( \|P\|_{\mathrm{op}} = 1 \).[web:87]
\end{proof}

\begin{definition}[Contraction]
A bounded operator \( T:\mathcal{H}\to\mathcal{H} \) is called a contraction if \( \|T\|_{\mathrm{op}} \le 1 \).[web:120]
\end{definition}

Every orthogonal projection is a contraction. More generally, any linear operator \( \Lambda_m:\mathcal{H} \to \mathcal{H} \) with \( \|\Lambda_m\|_{\mathrm{op}} \le 1 \) is a contraction and will be used as a \emph{spectral anchor} in the PITN dynamics.

\section{Norm Bounds for PITN Updates}

We now derive the norm bounds used in the main text to establish local stability of the PITN update rule under bounded operators and memory.

\subsection{Setup and assumptions}

Recall the PITN one-step update
\begin{equation}\label{eq:appendix-update}
T(t+1) = \Lambda_m\big(A_{p_{i_t}} T(t) + Q_{ent}(t)\big),
\end{equation}
where:
\begin{itemize}[leftmargin=*]
  \item \( T(t) \in \mathcal{H} \) is the state at time \( t \);
  \item \( A_{p_{i_t}} \in \mathcal{B}(\mathcal{H}) \) is the prime operator at step \( t \);
  \item \( \Lambda_m:\mathcal{H}\to\mathcal{H} \) is the spectral anchor (projection or contraction onto the variational manifold \( \mathcal{M} \subset \mathcal{H} \));
  \item \( Q_{ent}(t) \in \mathcal{H} \) is the ecological memory term.
\end{itemize}

We repeat the assumptions in a form suitable for analysis.

\begin{description}[leftmargin=2em]
\item[Assumption A1 (Bounded prime operators).]
There exists a constant \( c > 0 \) such that for all active primes \( p \in P_N(t) \) and all \( t \),
\[
\|A_p\|_{\mathrm{op}} \le c.
\]

\item[Assumption A2 (Contractive spectral anchor).]
The anchor \( \Lambda_m \) is a contraction:
\[
\|\Lambda_m\|_{\mathrm{op}} \le 1.
\]

\item[Assumption A3 (Uniformly bounded ecological memory).]
There exists a constant \( q_{\max} \ge 0 \) such that for all \( t \),
\[
\|Q_{ent}(t)\| \le q_{\max}.
\]
\end{description}

Assumption A1 is natural for bounded linear operators, and Assumption A2 holds in particular if \( \Lambda_m \) is an orthogonal projection onto a closed subspace \( \mathcal{M} \).[web:82][web:87] Assumption A3 encapsulates the effect of the finite history window and decaying kernel in the definition of \( Q_{ent} \).

\subsection{One-step norm inequality}

We first prove the basic one-step inequality used in Proposition 1 of the main text.

\begin{lemma}[One-step norm inequality]\label{lem:one-step}
Under Assumptions A1--A3, the PITN update \eqref{eq:appendix-update} satisfies
\[
\|T(t+1)\| \le c\,\|T(t)\| + q_{\max}.
\]
\end{lemma}

\begin{proof}
By Assumption A2,
\[
\|T(t+1)\|
= \|\Lambda_m(A_{p_{i_t}} T(t) + Q_{ent}(t))\|
\le \|\Lambda_m\|_{\mathrm{op}} \,\|A_{p_{i_t}} T(t) + Q_{ent}(t)\|
\le \|A_{p_{i_t}} T(t) + Q_{ent}(t)\|.
\]
Applying the triangle inequality and Assumption A1, we obtain
\[
\|T(t+1)\|
\le \|A_{p_{i_t}} T(t)\| + \|Q_{ent}(t)\|
\le \|A_{p_{i_t}}\|_{\mathrm{op}} \,\|T(t)\| + \|Q_{ent}(t)\|
\le c\,\|T(t)\| + q_{\max}.
\]
This is the desired inequality.
\end{proof}

\subsection{Affine recurrence and boundedness}

We now iterate the inequality from Lemma~\ref{lem:one-step} to derive uniform boundedness of the state norm when \( c < 1 \).

\begin{proposition}[Affine recurrence bound]\label{prop:affine-bound}
Let \( a_t = \|T(t)\| \). Under the conditions of Lemma~\ref{lem:one-step}, the sequence \( (a_t)_{t\ge 0} \) satisfies
\[
a_{t+1} \le c\,a_t + q_{\max}.
\]
If \( c < 1 \), then for all \( t \ge 0 \),
\[
a_t \le c^t a_0 + \frac{q_{\max}}{1-c} (1 - c^t),
\]
and in particular
\[
\limsup_{t\to\infty} a_t \le \frac{q_{\max}}{1-c}.
\]
\end{proposition}

\begin{proof}
The inequality \( a_{t+1} \le c\,a_t + q_{\max} \) is exactly Lemma~\ref{lem:one-step}. Solving this affine recurrence is standard: define
\[
b_t = a_t - \frac{q_{\max}}{1-c}.
\]
Then
\[
b_{t+1} = a_{t+1} - \frac{q_{\max}}{1-c}
\le c\,a_t + q_{\max} - \frac{q_{\max}}{1-c}
= c\,a_t + q_{\max} - \frac{q_{\max} - c q_{\max}}{1-c}
= c\,a_t - c\,\frac{q_{\max}}{1-c}
= c\,\left(a_t - \frac{q_{\max}}{1-c}\right)
= c\,b_t.
\]
Thus \( b_{t+1} \le c\,b_t \), implying by induction that
\[
b_t \le c^t b_0 = c^t\left(a_0 - \frac{q_{\max}}{1-c}\right).
\]
Therefore
\[
a_t 
= b_t + \frac{q_{\max}}{1-c}
\le c^t a_0 + \frac{q_{\max}}{1-c}(1 - c^t),
\]
and the claimed limit superior follows by taking \( t\to\infty \) with \( c^t \to 0 \) when \( c<1 \).
\end{proof}

As a direct corollary, we obtain the norm-stability result stated in the main body.

\begin{corollary}[Local norm stability]\label{cor:local-stability}
If \( c < 1 \), then the PITN dynamics defined by \eqref{eq:appendix-update} are norm-stable in the sense that
\[
\sup_{t\ge 0} \|T(t)\| < \infty
\]
and
\[
\|T(t)\| \to B \quad\text{with}\quad B \in \left[0,\,\frac{q_{\max}}{1-c}\right]
\]
in the sense that \( \limsup_{t\to\infty} \|T(t)\| \le \frac{q_{\max}}{1-c} \).
\end{corollary}

\begin{proof}
Immediate from Proposition~\ref{prop:affine-bound}.
\end{proof}

\subsection{Local linearization and spectral radius}

For completeness, we sketch the relation between operator-norm bounds and the spectral radius of a linearized dynamics around a fixed point, which underpins the ``spectral stability'' language in the main text.[web:87][web:121]

Let \( \bar T \in \mathcal{H} \) be a fixed point of the PITN dynamics (ignoring memory for the moment), i.e.\ it satisfies
\[
\bar T = \Lambda_m(A_{p^\ast} \bar T),
\]
for some fixed prime \( p^\ast \) or averaged operator.\footnote{In the presence of \( Q_{ent} \), one would first fix a stationary memory configuration or consider a reduced map.} Consider perturbations \( \delta T(t) = T(t) - \bar T \) and assume the update map is Fréchet differentiable at \( \bar T \). The linearization yields
\[
\delta T(t+1) \approx J\,\delta T(t),
\]
where \( J \in \mathcal{B}(\mathcal{H}) \) is the Jacobian. If we approximate
\[
J \approx \Lambda_m A_{p^\ast}
\]
and keep \( \Lambda_m \) contractive with \( \|\Lambda_m\|_{\mathrm{op}} \le 1 \), then
\[
\|J\|_{\mathrm{op}} \le \|\Lambda_m\|_{\mathrm{op}} \,\|A_{p^\ast}\|_{\mathrm{op}} \le c.
\]
Thus the spectral radius \( \rho(J) \), defined by
\[
\rho(J) := \sup\{|\lambda| : \lambda \in \sigma(J)\},
\]
satisfies \( \rho(J) \le \|J\|_{\mathrm{op}} \le c \).[web:87][web:121] Hence if \( c < 1 \), the linearized dynamics are locally asymptotically stable in the usual sense (perturbations decay in norm).

This provides a mathematically standard justification for speaking of \emph{spectral stability} or \emph{spectrally anchored} dynamics under the same operator-norm constraints used for the global norm bounds.

\section{Remarks on Non-Associativity and Path Dependence}

While composition of bounded operators in \( \mathcal{B}(\mathcal{H}) \) is associative by construction,[web:82][web:119] the PITN framework introduces non-associativity at the \emph{effective dynamics} level through two mechanisms:

\begin{enumerate}[leftmargin=*]
  \item Interleaving each application of a prime operator \( A_{p_{i_t}} \) with a projection/contraction \( \Lambda_m \) onto the variational manifold \( \mathcal{M} \), whose action need not commute with other operators.
  \item Coupling to the ecological memory term \( Q_{ent}(t) \), which depends explicitly on the history of prime words and inputs.
\end{enumerate}

Formally, for a given state \( T \), we compare
\[
T_1 = \Lambda_m\big(A_{p_2}\Lambda_m(A_{p_1} T + Q_{ent}^{(1)}) + Q_{ent}^{(2)}\big)
\]
and
\[
T_2 = \Lambda_m\big(A_{p_1}\Lambda_m(A_{p_2} T + Q_{ent}^{(1)}) + Q_{ent}^{(2)}\big),
\]
where \( Q_{ent}^{(1)}, Q_{ent}^{(2)} \) are history-dependent terms. In general, \( T_1 \neq T_2 \) even though \( A_{p_2}A_{p_1} = A_{p_2}A_{p_1} \) as formal products in \( \mathcal{B}(\mathcal{H}) \). This \emph{path dependence} is analogous to non-commuting update rules or non-ergodic dynamics in reinforcement learning, where the order of operations carries semantic information beyond the final state alone.[web:88]

\section{Summary of Bounds Used}

For rapid reference, we summarize the key operator-theoretic facts employed:

\begin{itemize}[leftmargin=*]
  \item \textbf{Operator norm.} For any \( A \in \mathcal{B}(\mathcal{H}) \), \( \|A\|_{\mathrm{op}} = \sup_{\|x\|=1} \|Ax\| \); this norm is submultiplicative: \( \|BA\|_{\mathrm{op}} \le \|B\|_{\mathrm{op}}\|A\|_{\mathrm{op}} \).[web:82][web:119]
  \item \textbf{Orthogonal projections.} Any orthogonal projection \( P \) on \( \mathcal{H} \) has \( \|P\|_{\mathrm{op}} = 1 \) and decomposes \( \mathcal{H} \) as \( \mathrm{ran}\,P \oplus \ker P \).[web:87]
  \item \textbf{Contractions.} Any operator with \( \|\Lambda_m\|_{\mathrm{op}} \le 1 \) is a contraction and does not increase vector norms.[web:120]
  \item \textbf{PITN update bound.} Under bounded prime operators and bounded memory, the PITN update obeys the affine recurrence \( \|T(t+1)\| \le c\,\|T(t)\| + q_{\max} \), yielding uniform boundedness and convergence to a ball of radius \( \frac{q_{\max}}{1-c} \) when \( c < 1 \).
  \item \textbf{Linearization and spectral radius.} For a Jacobian \( J \) satisfying \( \|J\|_{\mathrm{op}} \le c \), the spectral radius \( \rho(J) \le c \), implying local asymptotic stability when \( c<1 \).[web:87][web:121]
\end{itemize}

These results collectively justify the norm- and spectrum-based stability language in the PITN/QAGI framework while remaining within standard Hilbert-space operator theory.

\appendix
\section*{Mathematical Appendix}
\addcontentsline{toc}{section}{Mathematical Appendix}

\section{Operator-Theoretic Preliminaries}

In this appendix we collect and prove the basic operator-norm inequalities and stability bounds used in the main text. Throughout, \( \mathcal{H} \) denotes a complex Hilbert space with norm \( \|\cdot\| \) induced by its inner product, and \( \mathcal{B}(\mathcal{H}) \) denotes the Banach algebra of bounded linear operators on \( \mathcal{H} \) endowed with the operator norm.[web:87][web:119][web:124]

\subsection{Operator norm and submultiplicativity}

\begin{definition}[Operator norm]
Let \( A \in \mathcal{B}(\mathcal{H}) \). The operator norm of \( A \) is defined by
\[
\|A\|_{\mathrm{op}} := \sup_{\|x\| = 1} \|Ax\|.
\]
Equivalently, \( \|A\|_{\mathrm{op}} \) is the least constant \( c \ge 0 \) such that \( \|Ax\| \le c\|x\| \) for all \( x \in \mathcal{H} \).[web:82][web:119]
\end{definition}

\begin{lemma}[Submultiplicativity]\label{lem:submult}
If \( A,B \in \mathcal{B}(\mathcal{H}) \), then
\[
\|BA\|_{\mathrm{op}} \le \|B\|_{\mathrm{op}} \, \|A\|_{\mathrm{op}}.
\]
\end{lemma}

\begin{proof}
For any \( x \in \mathcal{H} \) with \( \|x\| = 1 \), we have
\[
\|BAx\| \le \|B\|_{\mathrm{op}} \, \|Ax\| \le \|B\|_{\mathrm{op}} \, \|A\|_{\mathrm{op}} \, \|x\|
= \|B\|_{\mathrm{op}} \, \|A\|_{\mathrm{op}}.
\]
Taking the supremum over all unit vectors \( x \) yields the result.[web:82][web:87]
\end{proof}

\subsection{Orthogonal projections and contractions}

We recall the basic facts about projections and contractions on Hilbert spaces.[web:87][web:120]

\begin{definition}[Projection]
A linear operator \( P:\mathcal{H} \to \mathcal{H} \) is a projection if \( P^2 = P \). It is an \emph{orthogonal projection} if, in addition, it is self-adjoint, i.e.\ \( P^\ast = P \).[web:87][web:119]
\end{definition}

\begin{theorem}[Projections and direct sum]\label{thm:projection}
If \( P:\mathcal{H}\to\mathcal{H} \) is an orthogonal projection, then
\[
\mathcal{H} = \mathrm{ran}\,P \oplus \ker P,
\]
and both \( \mathrm{ran}\,P \) and \( \ker P \) are closed subspaces.[web:87]
\end{theorem}

\begin{proposition}[Norm of an orthogonal projection]\label{prop:projection-norm}
If \( P \) is a nonzero orthogonal projection on \( \mathcal{H} \), then
\[
\|P\|_{\mathrm{op}} = 1.
\]
\end{proposition}

\begin{proof}
For any \( x \in \mathcal{H} \),
\[
\|Px\|^2 = \langle Px, Px \rangle = \langle x, P^\ast Px \rangle = \langle x, Px \rangle \le \|x\|\,\|Px\|
\]
by Cauchy--Schwarz, so \( \|Px\| \le \|x\| \), and hence \( \|P\|_{\mathrm{op}} \le 1 \).[web:87][web:122]
Since \( P \) is nonzero, there exists \( x \neq 0 \) with \( Px \neq 0 \). For such \( x \),
\[
\|Px\| = \|P(Px)\| \le \|P\|_{\mathrm{op}} \,\|Px\|.
\]
If \( \|P\|_{\mathrm{op}} < 1 \), then this implies \( \|Px\| = 0 \), a contradiction. Thus \( \|P\|_{\mathrm{op}} \ge 1 \). Combining with \( \|P\|_{\mathrm{op}} \le 1 \) yields \( \|P\|_{\mathrm{op}} = 1 \).[web:87]
\end{proof}

\begin{definition}[Contraction]
A bounded operator \( T:\mathcal{H}\to\mathcal{H} \) is called a contraction if \( \|T\|_{\mathrm{op}} \le 1 \).[web:120]
\end{definition}

Every orthogonal projection is a contraction. More generally, any linear operator \( \Lambda_m:\mathcal{H} \to \mathcal{H} \) with \( \|\Lambda_m\|_{\mathrm{op}} \le 1 \) is a contraction and will be used as a \emph{spectral anchor} in the PITN dynamics.

\section{Norm Bounds for PITN Updates}

We now derive the norm bounds used in the main text to establish local stability of the PITN update rule under bounded operators and memory.

\subsection{Setup and assumptions}

Recall the PITN one-step update
\begin{equation}\label{eq:appendix-update}
T(t+1) = \Lambda_m\big(A_{p_{i_t}} T(t) + Q_{ent}(t)\big),
\end{equation}
where:
\begin{itemize}[leftmargin=*]
  \item \( T(t) \in \mathcal{H} \) is the state at time \( t \);
  \item \( A_{p_{i_t}} \in \mathcal{B}(\mathcal{H}) \) is the prime operator at step \( t \);
  \item \( \Lambda_m:\mathcal{H}\to\mathcal{H} \) is the spectral anchor (projection or contraction onto the variational manifold \( \mathcal{M} \subset \mathcal{H} \));
  \item \( Q_{ent}(t) \in \mathcal{H} \) is the ecological memory term.
\end{itemize}

We repeat the assumptions in a form suitable for analysis.

\begin{description}[leftmargin=2em]
\item[Assumption A1 (Bounded prime operators).]
There exists a constant \( c > 0 \) such that for all active primes \( p \in P_N(t) \) and all \( t \),
\[
\|A_p\|_{\mathrm{op}} \le c.
\]

\item[Assumption A2 (Contractive spectral anchor).]
The anchor \( \Lambda_m \) is a contraction:
\[
\|\Lambda_m\|_{\mathrm{op}} \le 1.
\]

\item[Assumption A3 (Uniformly bounded ecological memory).]
There exists a constant \( q_{\max} \ge 0 \) such that for all \( t \),
\[
\|Q_{ent}(t)\| \le q_{\max}.
\]
\end{description}

Assumption A1 is natural for bounded linear operators, and Assumption A2 holds in particular if \( \Lambda_m \) is an orthogonal projection onto a closed subspace \( \mathcal{M} \).[web:82][web:87] Assumption A3 encapsulates the effect of the finite history window and decaying kernel in the definition of \( Q_{ent} \).

\subsection{One-step norm inequality}

We first prove the basic one-step inequality used in Proposition 1 of the main text.

\begin{lemma}[One-step norm inequality]\label{lem:one-step}
Under Assumptions A1--A3, the PITN update \eqref{eq:appendix-update} satisfies
\[
\|T(t+1)\| \le c\,\|T(t)\| + q_{\max}.
\]
\end{lemma}

\begin{proof}
By Assumption A2,
\[
\|T(t+1)\|
= \|\Lambda_m(A_{p_{i_t}} T(t) + Q_{ent}(t))\|
\le \|\Lambda_m\|_{\mathrm{op}} \,\|A_{p_{i_t}} T(t) + Q_{ent}(t)\|
\le \|A_{p_{i_t}} T(t) + Q_{ent}(t)\|.
\]
Applying the triangle inequality and Assumption A1, we obtain
\[
\|T(t+1)\|
\le \|A_{p_{i_t}} T(t)\| + \|Q_{ent}(t)\|
\le \|A_{p_{i_t}}\|_{\mathrm{op}} \,\|T(t)\| + \|Q_{ent}(t)\|
\le c\,\|T(t)\| + q_{\max}.
\]
This is the desired inequality.
\end{proof}

\subsection{Affine recurrence and boundedness}

We now iterate the inequality from Lemma~\ref{lem:one-step} to derive uniform boundedness of the state norm when \( c < 1 \).

\begin{proposition}[Affine recurrence bound]\label{prop:affine-bound}
Let \( a_t = \|T(t)\| \). Under the conditions of Lemma~\ref{lem:one-step}, the sequence \( (a_t)_{t\ge 0} \) satisfies
\[
a_{t+1} \le c\,a_t + q_{\max}.
\]
If \( c < 1 \), then for all \( t \ge 0 \),
\[
a_t \le c^t a_0 + \frac{q_{\max}}{1-c} (1 - c^t),
\]
and in particular
\[
\limsup_{t\to\infty} a_t \le \frac{q_{\max}}{1-c}.
\]
\end{proposition}

\begin{proof}
The inequality \( a_{t+1} \le c\,a_t + q_{\max} \) is exactly Lemma~\ref{lem:one-step}. Solving this affine recurrence is standard: define
\[
b_t = a_t - \frac{q_{\max}}{1-c}.
\]
Then
\[
b_{t+1} = a_{t+1} - \frac{q_{\max}}{1-c}
\le c\,a_t + q_{\max} - \frac{q_{\max}}{1-c}
= c\,a_t + q_{\max} - \frac{q_{\max} - c q_{\max}}{1-c}
= c\,a_t - c\,\frac{q_{\max}}{1-c}
= c\,\left(a_t - \frac{q_{\max}}{1-c}\right)
= c\,b_t.
\]
Thus \( b_{t+1} \le c\,b_t \), implying by induction that
\[
b_t \le c^t b_0 = c^t\left(a_0 - \frac{q_{\max}}{1-c}\right).
\]
Therefore
\[
a_t 
= b_t + \frac{q_{\max}}{1-c}
\le c^t a_0 + \frac{q_{\max}}{1-c}(1 - c^t),
\]
and the claimed limit superior follows by taking \( t\to\infty \) with \( c^t \to 0 \) when \( c<1 \).
\end{proof}

As a direct corollary, we obtain the norm-stability result stated in the main body.

\begin{corollary}[Local norm stability]\label{cor:local-stability}
If \( c < 1 \), then the PITN dynamics defined by \eqref{eq:appendix-update} are norm-stable in the sense that
\[
\sup_{t\ge 0} \|T(t)\| < \infty
\]
and
\[
\|T(t)\| \to B \quad\text{with}\quad B \in \left[0,\,\frac{q_{\max}}{1-c}\right]
\]
in the sense that \( \limsup_{t\to\infty} \|T(t)\| \le \frac{q_{\max}}{1-c} \).
\end{corollary}

\begin{proof}
Immediate from Proposition~\ref{prop:affine-bound}.
\end{proof}

\subsection{Local linearization and spectral radius}

For completeness, we sketch the relation between operator-norm bounds and the spectral radius of a linearized dynamics around a fixed point, which underpins the ``spectral stability'' language in the main text.[web:87][web:121]

Let \( \bar T \in \mathcal{H} \) be a fixed point of the PITN dynamics (ignoring memory for the moment), i.e.\ it satisfies
\[
\bar T = \Lambda_m(A_{p^\ast} \bar T),
\]
for some fixed prime \( p^\ast \) or averaged operator.\footnote{In the presence of \( Q_{ent} \), one would first fix a stationary memory configuration or consider a reduced map.} Consider perturbations \( \delta T(t) = T(t) - \bar T \) and assume the update map is Fréchet differentiable at \( \bar T \). The linearization yields
\[
\delta T(t+1) \approx J\,\delta T(t),
\]
where \( J \in \mathcal{B}(\mathcal{H}) \) is the Jacobian. If we approximate
\[
J \approx \Lambda_m A_{p^\ast}
\]
and keep \( \Lambda_m \) contractive with \( \|\Lambda_m\|_{\mathrm{op}} \le 1 \), then
\[
\|J\|_{\mathrm{op}} \le \|\Lambda_m\|_{\mathrm{op}} \,\|A_{p^\ast}\|_{\mathrm{op}} \le c.
\]
Thus the spectral radius \( \rho(J) \), defined by
\[
\rho(J) := \sup\{|\lambda| : \lambda \in \sigma(J)\},
\]
satisfies \( \rho(J) \le \|J\|_{\mathrm{op}} \le c \).[web:87][web:121] Hence if \( c < 1 \), the linearized dynamics are locally asymptotically stable in the usual sense (perturbations decay in norm).

This provides a mathematically standard justification for speaking of \emph{spectral stability} or \emph{spectrally anchored} dynamics under the same operator-norm constraints used for the global norm bounds.

\section{Remarks on Non-Associativity and Path Dependence}

While composition of bounded operators in \( \mathcal{B}(\mathcal{H}) \) is associative by construction,[web:82][web:119] the PITN framework introduces non-associativity at the \emph{effective dynamics} level through two mechanisms:

\begin{enumerate}[leftmargin=*]
  \item Interleaving each application of a prime operator \( A_{p_{i_t}} \) with a projection/contraction \( \Lambda_m \) onto the variational manifold \( \mathcal{M} \), whose action need not commute with other operators.
  \item Coupling to the ecological memory term \( Q_{ent}(t) \), which depends explicitly on the history of prime words and inputs.
\end{enumerate}

Formally, for a given state \( T \), we compare
\[
T_1 = \Lambda_m\big(A_{p_2}\Lambda_m(A_{p_1} T + Q_{ent}^{(1)}) + Q_{ent}^{(2)}\big)
\]
and
\[
T_2 = \Lambda_m\big(A_{p_1}\Lambda_m(A_{p_2} T + Q_{ent}^{(1)}) + Q_{ent}^{(2)}\big),
\]
where \( Q_{ent}^{(1)}, Q_{ent}^{(2)} \) are history-dependent terms. In general, \( T_1 \neq T_2 \) even though \( A_{p_2}A_{p_1} = A_{p_2}A_{p_1} \) as formal products in \( \mathcal{B}(\mathcal{H}) \). This \emph{path dependence} is analogous to non-commuting update rules or non-ergodic dynamics in reinforcement learning, where the order of operations carries semantic information beyond the final state alone.[web:88]

\subsection{Summary of Bounds Used}

For rapid reference, we summarize the key operator-theoretic facts employed:

\begin{itemize}[leftmargin=*]
  \item \textbf{Operator norm.} For any \( A \in \mathcal{B}(\mathcal{H}) \), \( \|A\|_{\mathrm{op}} = \sup_{\|x\|=1} \|Ax\| \); this norm is submultiplicative: \( \|BA\|_{\mathrm{op}} \le \|B\|_{\mathrm{op}}\|A\|_{\mathrm{op}} \).[web:82][web:119]
  \item \textbf{Orthogonal projections.} Any orthogonal projection \( P \) on \( \mathcal{H} \) has \( \|P\|_{\mathrm{op}} = 1 \) and decomposes \( \mathcal{H} \) as \( \mathrm{ran}\,P \oplus \ker P \).[web:87]
  \item \textbf{Contractions.} Any operator with \( \|\Lambda_m\|_{\mathrm{op}} \le 1 \) is a contraction and does not increase vector norms.[web:120]
  \item \textbf{PITN update bound.} Under bounded prime operators and bounded memory, the PITN update obeys the affine recurrence \( \|T(t+1)\| \le c\,\|T(t)\| + q_{\max} \), yielding uniform boundedness and convergence to a ball of radius \( \frac{q_{\max}}{1-c} \) when \( c < 1 \).
  \item \textbf{Linearization and spectral radius.} For a Jacobian \( J \) satisfying \( \|J\|_{\mathrm{op}} \le c \), the spectral radius \( \rho(J) \le c \), implying local asymptotic stability when \( c<1 \).[web:87][web:121]
\end{itemize}

These results collectively justify the norm- and spectrum-based stability language in the PITN/QAGI framework while remaining within standard Hilbert-space operator theory.

\section{Non-Associative Fractal Learning (NAFL)}
\addcontentsline{toc}{section}{Appendix B: Non-Associative Fractal Learning (NAFL)}

\subsection{Non-Associative Prime Operator Algebra}

We formalize Non-Associative Fractal Learning (NAFL) as a path-dependent dynamical system on a variational manifold \( \mathcal{M} \subset \mathcal{H} \), driven by a non-associative, prime-indexed operator grammar.

\begin{definition}[Prime operator alphabet]
Let \( \mathbb{P} \) denote the set of prime numbers. For each active prime \( p \in P_N(t) \subset \mathbb{P} \), define a bounded linear operator
\[
A_p \in \mathcal{B}(\mathcal{H}),
\]
where \( \mathcal{B}(\mathcal{H}) \) carries the operator norm \( \|\cdot\|_{\mathrm{op}} \) as before.[web:82][web:87]
\end{definition}

\begin{definition}[Prime words and composite operators]
A \emph{prime word} of length \( t \) is a sequence
\[
w(t) = (p_{i_1}, \dots, p_{i_t}), \qquad p_{i_k} \in \mathbb{P}.
\]
The associated composite operator in the underlying associative algebra \( \mathcal{B}(\mathcal{H}) \) is
\[
\mathcal{A}_{w(t)} := A_{p_{i_t}} A_{p_{i_{t-1}}} \cdots A_{p_{i_1}}.
\]
\end{definition}

At the level of \( \mathcal{B}(\mathcal{H}) \), operator multiplication is associative, i.e.
\[
(A_{p_2} A_{p_3}) A_{p_5} = A_{p_2} (A_{p_3} A_{p_5})
\]
as formal products.[web:87][web:151] NAFL introduces non-associativity at the level of \emph{effective dynamics} by interleaving each prime application with a stabilizing operator \( \Lambda_m \) and a history-dependent ecological term.

\begin{definition}[Stabilized prime operators]
Let \( \Lambda_m:\mathcal{H}\to\mathcal{H} \) be a contraction (spectral anchor) with \( \|\Lambda_m\|_{\mathrm{op}} \le 1 \). Define the stabilized prime operators
\[
\widetilde{A}_p := \Lambda_m A_p.
\]
\end{definition}

\begin{definition}[Effective NAFL update]
Given a state \( T \in \mathcal{M} \) and ecological memory \( Q_{ent} \), define a single NAFL step under prime \( p \) as
\[
\Phi_p(T) := \Lambda_m\big(A_p T + Q_{ent}\big).
\]
For a finite word \( w(t) = (p_{i_1}, \dots, p_{i_t}) \), define the effective NAFL composition
\[
\Phi_{w(t)} := \Phi_{p_{i_t}} \circ \cdots \circ \Phi_{p_{i_1}}.
\]
\end{definition}

\begin{definition}[Non-associativity at the dynamics level]
We say NAFL is \emph{dynamically non-associative} if, for some \( T \in \mathcal{M} \) and primes \( p_1,p_2,p_3 \),
\[
\Phi_{p_3}\big(\Phi_{p_2}(\Phi_{p_1}(T))\big) \neq
\Phi_{p_2}\big(\Phi_{p_3}(\Phi_{p_1}(T))\big),
\]
or more generally, if different bracketings and orderings of the same multiset of primes yield distinct effective compositions. 
\end{definition}

\begin{lemma}[Path dependence from non-associative updates]
Assume:
\begin{enumerate}[label=(\alph*), leftmargin=2em]
  \item \( \Lambda_m \) does not commute with all \( A_p \) (i.e.\ there exists \( p \) and \( T \) with \( \Lambda_m A_p T \neq A_p \Lambda_m T \));
  \item \( Q_{ent} \) depends nontrivially on the prime history (e.g.\ via a finite history window and nonzero decay kernel).
\end{enumerate}
Then there exist prime words \( w_1, w_2 \) with identical multisets of primes but different orderings such that
\[
\Phi_{w_1}(T) \neq \Phi_{w_2}(T),
\]
for some \( T \in \mathcal{M} \). In particular, the NAFL state at time \( t \) is intrinsically path-dependent.
\end{lemma}

\begin{proof}
By assumption (a), choose \( p,q \) and \( T \) such that \( \Lambda_m A_p T \neq A_p \Lambda_m T \). Consider length-two words \( w_1 = (p,q) \) and \( w_2 = (q,p) \). Ignoring memory for the moment, we have
\[
\Phi_{w_1}(T) = \Lambda_m\big(A_q \Lambda_m (A_p T)\big), \qquad
\Phi_{w_2}(T) = \Lambda_m\big(A_p \Lambda_m (A_q T)\big).
\]
If all operators commuted and associated trivially, these would coincide, but the non-commutation of \( \Lambda_m \) and \( A_p \) (and generically \( A_q \)) ensures that, for a suitable choice of \( T \), these expressions differ.

When \( Q_{ent} \) depends explicitly on the history, the asymmetry is stronger: the argument to \( Q_{ent} \) at each step depends on the sequence of primes applied up to that point, so swapping \( p \) and \( q \) changes both the operator nesting and the memory term. Thus there exist words with the same multiset but different orderings that yield distinct outputs, proving path dependence.
\end{proof}

In NAFL, the \emph{state descriptor} at time \( t \) is the pair \( (T(t), w(t)) \), where \( T(t) \in \mathcal{M} \) is the current tensor state and \( w(t) \) is the prime word applied so far. The same tensor \( T(t) \) with different histories \( w(t) \) can have different available continuations and different ecological couplings, reflecting the non-associative grammar.

\subsection{Fractal Tree of Prime Words}

We now formalize the ``fractal'' tree structure of prime words and its relation to the operator geometry.

\begin{definition}[Prime word tree]
Let \( \Sigma = P_N \subset \mathbb{P} \) be a finite active prime set. The \emph{prime word tree} \( \mathcal{T} \) is the rooted tree whose nodes are finite words \( w \in \Sigma^\ast \) (including the empty word \( \epsilon \)), with an edge from \( w \) to \( w' \) if \( w' = (w,p) \) for some \( p \in \Sigma \).
\end{definition}

Each node \( w \in \mathcal{T} \) is associated with a composite operator \( \mathcal{A}_w \) (in \( \mathcal{B}(\mathcal{H}) \)) and an effective NAFL map \( \Phi_w \) (on \( \mathcal{M} \)).

\begin{lemma}[Self-similar branching structure]
For any node \( w \in \mathcal{T} \), the subtree rooted at \( w \) is isomorphic, as a labeled tree, to the full tree \( \mathcal{T} \) (up to relabeling of the root), provided the active alphabet \( \Sigma \) is the same at all depths.
\end{lemma}

\begin{proof}
The subtree at \( w \) consists of all words of the form \( (w, u) \) with \( u \in \Sigma^\ast \). The mapping
\[
\varphi_w: \Sigma^\ast \to \{ \text{descendants of } w \}, \qquad u \mapsto (w,u),
\]
is a bijection that preserves adjacency: if \( u' = (u,p) \), then
\[
\varphi_w(u') = (w,u,p)
\]
is a child of \( \varphi_w(u) = (w,u) \). Thus the rooted subtree at \( w \) is isomorphic to the original tree rooted at \( \epsilon \).
\end{proof}

This combinatorial self-similarity underlies the fractal description: the same finite prime alphabet is used at all scales, and local subtrees resemble the whole. The geometry becomes nontrivial when we consider operator norms and stability.

\begin{definition}[Norm profile and spectral domain]
For each word \( w \), define its norm profile
\[
N(w) := \|\mathcal{A}_w\|_{\mathrm{op}},
\]
and its effective spectral domain as the set of states \( T \in \mathcal{M} \) such that the linearization of \( \Phi_w \) at \( T \) has spectral radius below a threshold (e.g.\ \( < 1 \)).
\end{definition}

The interaction between the branching combinatorics and the norm profile \( N(w) \) induces a fractal-like structure in the set of ``safe'' or ``good'' operator regions, with stability controlled by the boundedness of the \( A_p \) and the contractivity of \( \Lambda_m \) (see Appendix A) and the active alphabet controlled by a scheduling parameter \( \beta(t) \) as discussed in the main text.

\subsection{Decision-Theoretic Interpretation}

We now give a basic decision-theoretic formalization of NAFL in an RL-style setting, emphasizing the role of prime words as decision paths.

\begin{definition}[NAFL environment]
A NAFL environment is given by:
\begin{itemize}[leftmargin=*]
  \item A state space \( \mathcal{S} \), often taken as \( \mathcal{M} \times \Sigma^\ast \) (tensor state plus prime history).
  \item An observation space \( \mathcal{O} \) and observation function \( o_t = g(T(t), w(t)) \).
  \item A reward signal \( r_t \in \mathbb{R} \) and discount factor \( \gamma \in (0,1) \).
\end{itemize}
\end{definition}

\begin{definition}[Prime-selecting policy]
A \emph{prime-selecting policy} is a stochastic kernel
\[
\pi(p \mid o_t, w(t))
\]
giving the probability of choosing prime \( p \in P_N(t) \) conditioned on the current observation \( o_t \) and prime word \( w(t) \). A trajectory is a sequence
\[
\big(T(0), w(0) = \epsilon\big) \xrightarrow{p_{i_1}} \big(T(1), w(1)\big) \xrightarrow{p_{i_2}} \cdots
\]
with \( T(t+1) = \Phi_{p_{i_t}}(T(t)) \) and \( w(t+1) = (w(t), p_{i_t}) \).
\end{definition}

\begin{definition}[Value function over prime words]
The NAFL value function at time \( t \) can be written as
\[
V(w(t)) := \mathbb{E}\!\left[ \sum_{k \ge t} \gamma^{k-t} r_k \,\middle|\, \mathcal{A}_{w(t)}, \pi \right],
\]
with expectation over environment transitions and future prime selections under policy \( \pi \).
\end{definition}

Crucially, \( w(t) \) appears explicitly in the conditioning: two trajectories that lead to the same tensor state \( T(t) \) but arise from different words \( w(t) \) may have different value functions because the set of admissible continuation operators, ecological couplings, or safety constraints may differ.

\subsection{NAFL vs.\ Classical RL}

For contrast, classical RL typically treats the state as an element of \( \mathcal{S} \) (possibly with a recurrent hidden state), with an underlying \emph{associative} composition of updates at the parameter level (e.g.\ weight matrices in a neural network).[web:152][web:155] NAFL differs in two key ways:

\begin{itemize}[leftmargin=*]
  \item \textbf{Explicit non-associativity and path memory.} NAFL does not collapse different action histories into a single effective operator; the prime-word tree preserves distinct branches indefinitely, with non-associative effective dynamics ensuring that order and bracketing matter.
  \item \textbf{Hierarchical structure as intrinsic branching.} Hierarchical policies emerge from the tree structure of prime words, where repeated subwords act as reusable skills or motifs, rather than from explicitly hand-crafted options or macro-actions.[web:139][web:152][web:157]
\end{itemize}

This framing aligns with existing work on non-associative algebras[web:147][web:148][web:151] and path-dependent learning architectures[web:130][web:145], while embedding them in a tensor-network-based RL substrate.

\subsection{Implementable Structures (Sketch)}

To connect NAFL to implementable tensor-network learning:

\begin{itemize}[leftmargin=*]
  \item Each prime operator \( A_p \) can be realized as a small tensor-network module (e.g.\ MPS/TTN block) with prime-specific parameters, building on supervised learning with tensor networks.[web:150][web:153][web:156]
  \item The manifold \( \mathcal{M} \) can be instantiated as a class of tensor-network states with fixed bond dimension, with updates carried out via TDVP-style or gradient-based optimization.[web:51][web:76][web:156]
  \item The policy \( \pi \) can be implemented as a neural network taking as input an encoding of the current word \( w(t) \) (e.g.\ via a tree encoder or sequence model) plus observation features, and outputting a distribution over primes in \( P_N(t) \).[web:152][web:143]
  \item Ecological coupling \( Q_{ent} \) can be implemented as a learned symmetric kernel over prime words, penalizing or encouraging joint updates to related branches based on structural similarity or shared reward patterns.
\end{itemize}

This preserves the theoretical pillars—prime indexing, non-associativity, spectral anchoring via \( \Lambda_m \), and fractal branching—while mapping them onto concrete architectures informed by existing tensor-network ML and hierarchical RL practice.[web:150][web:152][web:156]



\end{document}